\section{Visualisasi Data}
\label{sec:visualisasidata}

Visualisasi data merupakan representasi grafis dari data untuk membantu orang memahami konteks dan signifikansi. Tujuan visualisasi data sendiri yakni untuk menyampaikan informasi secara ringkas dan jelas sehingga mudah dipahami oleh pembaca \cite{alghivary2023visualisasi}. Dengan menggunakan elemen visual seperti bagan, grafik, dan peta, visualisasi data menyediakan cara yang mudah diakses untuk melihat dan memahami tren, pencilan (outliers), dan pola dalam data. Dalam lingkungan organisasi, visualisasi data memungkinkan para pengambil keputusan untuk melihat analitik yang disajikan secara visual, sehingga mereka dapat memahami konsep yang sulit atau mengidentifikasi pola baru.

Prinsip visualisasi data diterapkan secara langsung pada User Interface (UI) website MonPD Reborn. Seluruh halaman pada website ini menggunakan berbagai komponen visual seperti card, tabel, grafik, dan lainnya untuk menampilkan data statistik, diagram batang untuk perbandingan, dan tabel data berwarna untuk status. Tujuan utamanya adalah agar pengguna dapat dengan cepat memahami kondisi terkini dari realisasi dan pemeriksaan pajak tanpa harus membaca laporan tekstual yang panjang dan kompleks.
